\section{Deadly Triad 让自举误差在分布外放大}
\label{sec:16-Control-and-Reinforcement-Learning/02-Reinforcement-Learning/06}

一份离线日志摆在面前：三千万条转移，全部由半年前上线的那版策略产生。你要在这批数据上训一个 Q 网络，用来评估一个改过的新策略。前两万步一切正常，TD 损失稳稳下降；第五万步左右，网络输出的 Q 值从个位数爬到几百，第八万步爬过 \(10^{5}\)，而 TD 损失看上去反倒更小了。代码里没有除零，梯度范数也不算大，换随机种子只把崩溃点推后几万步；学习率砍到十分之一，崩得慢了，方向没变。

这里没有 bug。发散来自三样东西的组合，而这三样都是你为了让问题可解才请进来的。

\subsection{任意两个都安全，三个凑齐才危险}

三个角是这样的。\textbf{函数逼近}让你用一组参数 \(w\) 表示整个状态空间上的价值，代价是不同状态共用参数。\textbf{自举}让你把目标写成 \(r+\gamma\widehat v_w(s')\)，奖励信息不必等到回合结束，代价是目标里含有你自己的当前估计。\textbf{离策略}让你复用旧策略、他人策略甚至人工日志产生的数据，代价是更新时看到的状态分布来自行为策略 \(\mu\)，而你想要的答案属于目标策略 \(\pi\)。

值得记住的结构是：这三样东西，任意取两样，都有像样的收敛结论；三样凑齐，期望更新本身就可能没有吸引不动点。

只要 \(\pi=\mu\)，线性函数逼近加自举给出 on-policy TD，上一节已经证过它收敛到投影 Bellman 方程的解。只要放弃自举，改用 Monte Carlo 回报做目标，离策略下的函数逼近退化成一个带权重的最小二乘回归，目标不随参数动，收敛性来自随机逼近的常规条件。只要放弃函数逼近改用表格，每个状态一个独立参数，离策略 Q 学习在标准步长条件下以概率一收敛。三条边各自安全，三角形的内部危险。

\begin{center}
\begin{tikzpicture}[>=stealth,scale=0.9]
\coordinate (FA) at (0,2.4);
\coordinate (BS) at (-2.2,-1.1);
\coordinate (OP) at (2.2,-1.1);
\draw[thick] (FA) -- (BS) -- (OP) -- cycle;
\filldraw (FA) circle (2.2pt);
\filldraw (BS) circle (2.2pt);
\filldraw (OP) circle (2.2pt);
\node[above] at (FA) {函数逼近};
\node[below] at (BS) {自举};
\node[below] at (OP) {离策略};
\node[left,font=\small,align=right] at (-1.45,0.6) {on-policy\\线性 TD};
\node[right,font=\small,align=left] at (1.45,0.6) {Monte Carlo\\回归};
\node[font=\small] at (0,-1.75) {表格 Q 学习};
\node[font=\small,align=center] at (0,0.35) {三角齐全\\可能发散};
\begin{scope}[shift={(6.4,-1.1)}]
\draw[->] (0,0) -- (5.0,0) node[below,font=\small] {迭代 \(k\)};
\draw[->] (0,0) -- (0,3.5) node[left,font=\small] {\(\abs{w_k}\)};
\draw[thick] plot[domain=0:4.4,samples=45] (\x,{0.12*exp(0.72*\x)});
\draw[dashed] plot[domain=0:4.4,samples=45] (\x,{1.4*exp(-0.9*\x)});
\node[font=\small] at (2.1,2.7) {三角齐全};
\node[font=\small] at (3.4,0.42) {去掉任一角};
\end{scope}
\end{tikzpicture}
\end{center}

\emph{三条边各自有收敛结论，内部那一点没有；去掉任一角，误差按虚线回落，三角齐全时按实线指数放大。}

图上三条边的名字都是熟悉的算法，这一点值得多看一眼。Deadly Triad 说的不是某个算法有缺陷，它说的是三种便利叠加之后有一条保证消失了。你想要大状态空间、想要每步就能学、又想复用别人的数据，代价写在这张图的中心。

\subsection{误差沿一条闭合回路被放大}

三个角各自贡献一段链路，合起来闭成一个圈。

自举把目标接到自己身上：\(s\) 的更新目标里含 \(\widehat v_w(s')\)，抬高 \(s'\) 的估计就抬高 \(s\) 的目标。函数逼近让这种抬高外溢：\(s\) 与 \(s'\) 共享参数，为拟合 \(s\) 而调 \(w\)，\(s'\) 的估计跟着动，误差在状态之间串扰。离策略决定了谁被抬、谁被拟合：投影按 \(\mu\) 的访问频率加权，权重大的状态被反复拟合，权重小的状态几乎不被纠正。

于是可以出现这样一圈：更新为了拟合 \(\mu\) 常访问的状态而调参数，副作用是把 \(\mu\) 极少访问、\(\pi\) 却常去的那个后继状态的估计抬高；抬高的估计通过自举回流到目标里，让 \(\mu\) 常访问的状态需要再抬一次。分布外的那个状态永远等不来一次采样纠正，误差每绕一圈乘上一个大于一的系数。缺哪一角，圈都闭不上：没有自举，抬高不会回流；没有函数逼近，抬高不会外溢；没有分布错配，被抬高的状态自己也会被采样纠正。

\subsection{两个状态就足以发散}

不需要神经网络。取两个状态、零奖励、\(\gamma=0.9\)，目标策略的转移矩阵是

\[
P_\pi=
\begin{bmatrix}
0&1\\
0&1
\end{bmatrix},
\]

即无论现在在哪，下一步都到状态 2。真实价值恒为零。只用一个特征，

\[
\phi=\begin{bmatrix}1\\2\end{bmatrix},
\qquad
\widehat v(w)=\phi w,
\]

零显然可表示：\(w=0\) 即可。数据由另一行为策略产生，长期访问权重为 \(D_\mu=\diag(0.99,0.01)\)，也就是说日志几乎全落在状态 1 上，而目标策略几乎全在状态 2。

把 \(T_\pi(\phi w)=\gamma P_\pi\phi w\) 沿 \(D_\mu\) 加权投影回特征直线，参数上的迭代是一个纯粹的乘法：

\[
w_{k+1}=cw_k,
\qquad
c=\gamma\,\frac{\phi^{\trans}D_\mu P_\pi\phi}{\phi^{\trans}D_\mu\phi}
=0.9\times\frac{2.02}{1.03}
\approx1.765.
\]

唯一不动点是 \(0\)，可从任何 \(w_0\neq0\) 出发，误差每步乘 \(1.765\)。写成期望 TD 更新也一样：\(A=\phi^{\trans}D_\mu(I-\gamma P_\pi)\phi=1.03-1.818=-0.788\)，平均更新 \(w\leftarrow w+\alpha(0-Aw)=(1+0.788\alpha)w\) 稳定地朝远离零的方向走。步长只改变发散速度。

同一个例子可以逐个拆角验证前面那张图。把权重换成目标策略自己的平稳分布 \(D_\pi=\diag(0,1)\)，特征不变、自举不变，系数变成 \(0.9\times 4/4=0.9<1\)，迭代收敛；发散不在特征上，也不在自举上。保留离策略数据但把目标换成 Monte Carlo 回报，目标恒为零且不含 \(w\)，回归立刻把 \(w\) 拉向零。改用表格表示，两个状态各有一个参数，状态 2 的参数被自己的样本压回零，状态 1 不再被牵连。三次拆解，三次收敛。

\subsection{压缩性丢在了范数不匹配上}

值迭代能收敛，靠的是 \(T_\pi\) 在 sup 范数下是 \(\gamma\)-压缩：\(\norm{T_\pi v-T_\pi u}_\infty\le\gamma\norm{v-u}_\infty\)。函数逼近下实际反复作用的算子是 \(\Pi_D T_\pi\)，投影 \(\Pi_D\) 在自己的 \(D\)-加权范数下不扩张。两处保证写在两种几何里，复合起来未必压缩。

补上一条就够了：只要 \(P_\pi\) 在 \(D\)-范数下不扩张，复合算子就是 \(\gamma\)-压缩。骨架是两步。由 Jensen 不等式，

\[
\begin{aligned}
\norm{P_\pi v}_{d}^{2}
&=\sum_s d(s)\Big(\sum_{s'}P_\pi(s'\given s)\,v(s')\Big)^{2}\\
&\le\sum_s d(s)\sum_{s'}P_\pi(s'\given s)\,v(s')^{2}
=\sum_{s'}(dP_\pi)(s')\,v(s')^{2}.
\end{aligned}
\]

若 \(dP_\pi=d\)，右端恰是 \(\norm{v}_d^{2}\)，不扩张成立。条件对账：第一步只用了平方的凸性，任何 \(d\) 都成立；第二步要求 \(d\) 是 \(P_\pi\) 的平稳分布，只有 on-policy 才自动满足。离策略时右端可以大出一个因子 \(\max_s (d_\mu P_\pi)(s)/d_\mu(s)\)，在上面的例子里 \(d_\mu P_\pi=(0,1)\)，这个比值是 \(1/0.01=100\)，\(\gamma\) 的收缩根本不够抵。

这也解释了为什么覆盖不足格外危险：\(d_\mu\) 在 \(\pi\) 常去的状态上越接近零，这个比值越大。行为策略与目标策略差多远，是可以量化的，后面讲 off-policy 校正的一节会把它写成重要性权重的方差。

\subsection{发散、过估计与震荡要分开记账}

Q 值变大不一定是 Deadly Triad。对含噪估计取最大值会产生系统性的乐观偏差，这是 \(\E[\max]\ge\max\E\) 的老问题，与自举无关；固定步长会让参数在稳定点附近持续抖动，这在 \hyperref[sec:09-Convex-Optimization/08-Nonconvex-and-Stochastic-Optimization/04]{``SGD 在期望中下降但噪声会留下平台''} 里讨论过；奖励尺度过大、网络初始化不当造成的数值溢出更是另一码事。Deadly Triad 说的是期望更新的方向本身出了问题，把采样噪声降到零也修不好。

分辨它们靠几个可观察的信号：

\begin{itemize}
\tightlist
\item 把 batch 增大十倍，方差明显下降而均值轨迹照旧远离合理范围，问题在平均动力学而非噪声；
\item Q 值远超环境可能给出的最大折扣回报，且 TD 目标与预测同步上抬，说明自举正在自我抬升；
\item 把数据换成当前策略采集的 on-policy 样本后立刻稳定，说明分布错配是关键那一角。
\end{itemize}

最后一条尤其值得跑一次。它花的成本不高，却能把``理论上的三角''和``工程上的脏数据''区分开来。只有个别种子崩溃，则更该先查探索覆盖与数据清洗。

\subsection{止损的手段各自动了不同的角}

按角来分类，现有办法就清楚了。Gradient-TD、TDC、GTD2 把目标从``TD 更新''换成一个真有梯度的目标函数（投影 Bellman 误差），于是随机梯度下降的常规结论重新可用；emphatic TD 走另一条路，用一组强调权重把采样分布重新加权，人为恢复上面那条平稳性条件。残差梯度法直接对 Bellman 误差求梯度，语义干净，代价是需要同一状态的两次独立后继采样，方差大、收敛慢。把 \(\gamma\) 调小则是缩短有效视界，自举链条短了，放大系数也小了，代价是你换了一个问题在解。

目标网络和经验回放不属于以上任何一类。它们不恢复任何保证，只是把自举目标的移动速度和样本的相关性压下去，让崩溃的时间尺度慢于训练收敛的时间尺度。这是工程缓解，不是定理。下一节的 DQN 把这两件事做到了当时的极致，读它的时候值得始终带着一个问题：被放慢的到底是哪一段耦合。
